\cotefile{97}
\scanpage{5}

Classes de Chern de la représentation régulière et ses multiples.

$G$ groupe cyclique d'ordre $n$,
\[
H^*(G,\mathbb Z)^+ \simeq
\underline{\mathrm{Sym}}_{\mathbb Z_{\! /n\mathbb Z}}
  (\overset{\text{\scriptsize\inserted{dual sur $\mathbb Z_{\! /n\mathbb Z}$}}}{\check G})^+
\simeq \mathbb Z_{\! /n\mathbb Z}[\xi]^+
\]
où $\xi$ est un générateur de
$H^2(G,\mathbb Z)\simeq
\underset{\substack{\rotatebox{90}{$\simeq$}\\ \check G}}{
\mathrm H^{\deleted{\unsure{2|*}}}\mathrm{om}(G,\mathbb Q/\mathbb Z)}$

La représentation régulière \inserted{(sur $\mathbb C$)} \deleted{est} correspond à la somme des représentations de degré $1$ \unsure{associés|associées} aux \unsure{différents} $\xi\in\check G$.

On a donc
\[
c(r_G)=\prod_{\xi\in\check G}(1+\xi)
=\prod_{\substack{\deleted{{\unsure{k\in\mathbb Z|k=0}}^{\unsure{n-1|n}}}\\ \lambda\in\mathbb Z_{\! /n\mathbb Z}}}(1+\lambda\xi_0)
\deleted{=\unsure{\sum_{i=0}^{n}c_i\xi_0^i\;(\text{dans }\mathbb Z_{\! /n\mathbb Z}[\xi])}}
\]

Supposons $n=p$ premier. On trouve
\[
\left\{\begin{aligned}
c_p(r_G)&=0\\
c_{p-1}(r_G)&=\prod_{\xi\in\check G-\{0\}}\xi
=\left(\prod_{\lambda\in\mathbb F_p^{\unsure{+|*}}}\lambda\right)\xi_0^{p-1}
=\xi_0^{p-1}\quad \text{(indépendant du choix de $\xi_0\in\check G$ !)}
\end{aligned}\right.
\]

D'où
\[
\left\{\begin{aligned}
c_{\deleted{\unsure{p}}i}(Nr_G)&=0\quad\text{si }i>N(p-1)\\
c_{N(p-1)}(Nr_G)&=\xi^{N(p-1)}\qquad(\xi\in\check G)
\end{aligned}\right.
\]

Soit maintenant $G$ un groupe fini\inserted{ d'ordre $n$}, d'ordre divisible par $p$, \deleted{Alors on \unsure{trouve|tous}} et soit $H$ un sous-groupe d'ordre $p$. Considérons $n'=n/p$, alors $\deleted{(Nr_G)\vert}r_G\vert H=n'r_H$ donc $(Nr_G)\vert H=(Nn')r_H$, d'où $\bigl(c_i(nr_G)\bigr)\vert H=c_i(Nn'r_H)$, et par suite
\[
c_{Nn'(p-1)}(r_G)\vert H\ne0,\quad\text{à fortiori}\quad c_{Nn'(p-1)}(r_G)\ne0
\]
Donc
\[
\boxed{H^{\deleted{\unsure{*|2}}N\unsure{K|k}}(G)(p)\ne0\quad\text{pour tout entier }N,\quad\text{où}\quad k=2n'(p-1)}
\]
% archive sheet number: 4
